\documentclass{article}
\usepackage{amsmath,amssymb,amsthm}
\usepackage{enumitem}
\usepackage{hyperref}
\usepackage{graphicx}
\usepackage{bm}

\newtheorem{theorem}{Theorem}
\newthetheorem{lemma}{Lemma}
\newtheorem{assumption}{Assumption}
\newtheorem{corollary}{Corollary}

\begin{document}

%%%%%%%%%%%%%%%%%%%%%%%%%%%%%%%%%%%%%%%%%%%%%%%%%%%%%%%%%%%%%%%%%%%%%%%%%%%%%
\section*{Algorithm Name}
\textbf{Accelerated Proximal Gradient with Momentum (APG-M)}  
(also known as Nesterov’s Accelerated Proximal Gradient).

%%%%%%%%%%%%%%%%%%%%%%%%%%%%%%%%%%%%%%%%%%%%%%%%%%%%%%%%%%%%%%%%%%%%%%%%%%%%%
\section*{Mathematical Formulation}
Let  

\[
F(x)=f(x)+g(x),\qquad f:\mathbb R^{d}\to\mathbb R \text{ convex, $L$–smooth},
\qquad g:\mathbb R^{d}\to\mathbb R\cup\{+\infty\}\text{ convex (possibly non‑smooth)} .
\]

Denote by  

\[
\operatorname{prox}_{\eta g}(v)=\arg\min_{u\in\mathbb R^{d}}\Big\{ g(u)+\frac{1}{2\eta}\|u-v\|^{2}\Big\}
\]

the proximal operator of $g$ with stepsize $\eta>0$.  
The algorithm generates two sequences $\{x_{k}\}$ and $\{y_{k}\}$:

\[
\boxed{
\begin{aligned}
y_{k}      &= x_{k}+ \beta_{k}\,(x_{k}-x_{k-1}), \qquad (k\ge 0)\\[2mm]
x_{k+1} &= \operatorname{prox}_{\eta g}\!\bigl(y_{k}-\eta \nabla f(y_{k})\bigr),
\end{aligned}}
\tag{APG-M}
\]

where $x_{-1}=x_{0}$ and $\beta_{k}\in[0,1)$ is a momentum coefficient that will be
chosen according to a sequence $\{t_{k}\}$ (see below).

%%%%%%%%%%%%%%%%%%%%%%%%%%%%%%%%%%%%%%%%%%%%%%%%%%%%%%%%%%%%%%%%%%%%%%%%%%%%%
\section*{Pseudocode}
\begin{enumerate}[label=Step \arabic*:, leftmargin=*, align=left]
\item \textbf{Input:} initial point $x_{0}\in\mathbb R^{d}$, stepsize $\eta\in(0,1/L]$, 
      integer $K$ (max. iterations).
\item Set $x_{-1}\leftarrow x_{0}$, $t_{0}\leftarrow 1$.
\item \textbf{For} $k=0,1,\dots ,K-1$ \textbf{do}
\begin{enumerate}[label=\arabic*.]
\item $t_{k+1}\leftarrow\frac{1+\sqrt{1+4t_{k}^{2}}}{2}$,
\item $\beta_{k}\leftarrow\frac{t_{k}-1}{t_{k+1}}$,
\item $y_{k}\leftarrow x_{k}+ \beta_{k}(x_{k}-x_{k-1})$,
\item $x_{k+1}\leftarrow\operatorname{prox}_{\eta g}\!\bigl(y_{k}-\eta\nabla f(y_{k})\bigr)$.
\end{enumerate}
\item \textbf{Output:} $x_{K}$ (or the best iterate in terms of $F$).
\end{enumerate}

%%%%%%%%%%%%%%%%%%%%%%%%%%%%%%%%%%%%%%%%%%%%%%%%%%%%%%%%%%%%%%%%%%%%%%%%%%%%%
\section*{Dry Run Example}
Consider the scalar problem  

\[
f(w)=(w-3)^{2},\qquad g\equiv0,
\]
so $\operatorname{prox}_{\eta g}(v)=v$.  
Here $L=2$ (since $\nabla f(w)=2(w-3)$ is $2$‑Lipschitz).  
Choose $\eta=0.1\le 1/L$, $x_{0}=0$, and run three iterations with the
standard Nesterov schedule.

\begin{center}
\begin{tabular}{c|c|c|c|c}
$k$ & $t_{k}$ & $\beta_{k}$ & $y_{k}$ & $x_{k+1}=y_{k}-\eta\nabla f(y_{k})$ \\ \hline
0 & $1$ & $0$ & $0$ & $0-0.1\cdot 2(0-3)=0.6$ \\
1 & $1.618$ & $\frac{1-1}{1.618}=0$ & $0.6$ & $0.6-0.1\cdot2(0.6-3)=0.6-0.1\cdot(-4.8)=1.08$ \\
2 & $2.193$ & $\frac{1.618-1}{2.193}=0.281$ & $1.08+0.281(1.08-0.6)=1.08+0.281\cdot0.48\approx1.215$ &
$1.215-0.1\cdot2(1.215-3)=1.215-0.1\cdot(-3.57)=1.572$ 
\end{tabular}
\end{center}

Objective values $F(x_{k})=f(x_{k})$:

\[
f(0)=9,\quad f(0.6)=5.76,\quad f(1.08)=3.6864,\quad f(1.572)=2.040\; .
\]

The values decrease rapidly, illustrating the acceleration effect.

%%%%%%%%%%%%%%%%%%%%%%%%%%%%%%%%%%%%%%%%%%%%%%%%%%%%%%%%%%%%%%%%%%%%%%%%%%%%%
\section*{Assumptions}
\begin{assumption}[Smoothness of $f$]\label{as:smooth}
$f$ is convex and differentiable with $L$‑Lipschitz continuous gradient, i.e.  

\[
\|\nabla f(x)-\nabla f(y)\|\le L\|x-y\|\qquad\forall x,y\in\mathbb R^{d}.
\tag{A1}
\]
\end{assumption}

\begin{assumption}[Convexity of $g$]\label{as:convex}
$g$ is proper, closed and convex (possibly non‑smooth).  
Its proximal operator \eqref{eq:prox} is well defined for any $\eta>0$.
\tag{A2}
\end{assumption}

\begin{assumption}[Stepsize]\label{as:stepsize}
\[
0<\eta\le\frac{1}{L}.
\tag{A3}
\]
\end{assumption}

\begin{assumption}[Momentum coefficients]\label{as:beta}
Define a scalar sequence $\{t_{k}\}_{k\ge0}$ by  

\[
t_{0}=1,\qquad t_{k+1}= \frac{1+\sqrt{1+4t_{k}^{2}}}{2},
\tag{A4}
\]
and set  

\[
\beta_{k}= \frac{t_{k}-1}{t_{k+1}}\in[0,1).
\tag{A5}
\]
\end{assumption}

All subsequent results rely exclusively on \ref{as:smooth}--\ref{as:beta}.

%%%%%%%%%%%%%%%%%%%%%%%%%%%%%%%%%%%%%%%%%%%%%%%%%%%%%%%%%%%%%%%%%%%%%%%%%%%%%
\section*{Main Convergence Theorem}
\begin{theorem}[Accelerated rate]\label{thm:rate}
Let $\{x_{k}\}$ be generated by \eqref{APG-M} under Assumptions
\ref{as:smooth}–\ref{as:beta}.  
Denote $x^{\star}\in\arg\min_{x}F(x)$ and $D:=\|x_{0}-x^{\star}\|$.  
Then for every $k\ge1$  

\[
F(x_{k})-F(x^{\star})\;\le\;\frac{2L D^{2}}{(k+1)^{2}} .
\tag{T1}
\]

Consequently $F(x_{k})\to F(x^{\star})$ with the optimal $O(1/k^{2})$ rate for
first‑order methods on composite convex problems.
\end{theorem}

\begin{corollary}[Ergodic bound]
Define the weighted average $\bar x_{K}= \frac{2}{K(K+1)}\sum_{k=1}^{K}k\,x_{k}$.  
Then  

\[
F(\bar x_{K})-F(x^{\star})\le \frac{2L D^{2}}{K^{2}} .
\tag{C1}
\]
\end{corollary}

%%%%%%%%%%%%%%%%%%%%%%%%%%%%%%%%%%%%%%%%%%%%%%%%%%%%%%%%%%%%%%%%%%%%%%%%%%%%%
\section*{Theoretical Properties}
\begin{itemize}
\item \textbf{Stability.} The recursion \eqref{APG-M} is a non‑expansive mapping
      in the metric induced by the Bregman distance associated with $f$,
      guaranteeing that the iterates remain bounded (by $D$).

\item \textbf{Hyper‑parameter sensitivity.}  
      The stepsize bound $\eta\le1/L$ is tight: any larger $\eta$ may violate the
      descent Lemma (Lemma~\ref{lem:descent}) and break the $O(1/k^{2})$ rate.  
      The momentum schedule \eqref{A4}–\eqref{A5} is the unique choice that
      yields the exact telescoping structure used in the proof; other
      constant $\beta$ give only $O(1/k)$.

\item \textbf{Comparison to baselines.}  
      \begin{enumerate}
        \item \textbf{Plain proximal gradient (PG).}  PG satisfies
              $F(x_{k})-F(x^{\star})\le \frac{L D^{2}}{2k}$, i.e. a slower $1/k$ rate.
        \item \textbf{Accelerated gradient without prox (Nesterov).}  Theorem~\ref{thm:rate}
              recovers the classic $O(1/k^{2})$ bound when $g\equiv0$.
        \item \textbf{Heavy‑ball momentum.}  Does not guarantee the same worst‑case
              rate for composite problems; APG-M is provably optimal in the
              first‑order oracle model.
      \end{enumerate}
\end{itemize}

%%%%%%%%%%%%%%%%%%%%%%%%%%%%%%%%%%%%%%%%%%%%%%%%%%%%%%%%%%%%%%%%%%%%%%%%%%%%%
\section*{Preliminaries (Definitions and Lemmas)}
\begin{lemma}[Descent Lemma]\label{lem:descent}
If $f$ satisfies (A1) and $\eta\le 1/L$, then for any $u,v\in\mathbb R^{d}$  

\[
f(u)\le f(v)+\langle\nabla f(v),u-v\rangle+\frac{L}{2}\|u-v\|^{2}
\le f(v)+\langle\nabla f(v),u-v\rangle+\frac{1}{2\eta}\|u-v\|^{2}.
\tag{L1}
\]
\end{lemma}
\begin{proof}
The first inequality is classical; the second follows from $\eta\le1/L$,
i.e. $L\le 1/\eta$.
\end{proof}

\begin{lemma}[Three‑Point Property of the Proximal Operator]\label{lem:prox}
For any $v\in\mathbb R^{d}$ and any $u\in\mathbb R^{d}$, letting
$u^{+}= \operatorname{prox}_{\eta g}(v)$ we have  

\[
g(u^{+})+ \frac{1}{2\eta}\|u^{+}-v\|^{2}\le
g(u)+\frac{1}{2\eta}\|u-v\|^{2}
-\frac{1}{2\eta}\|u-u^{+}\|^{2}.
\tag{L2}
\]
\end{lemma}
\begin{proof}
Standard optimality condition of the proximal subproblem.
\end{proof}

\begin{lemma}[Potential Function]\label{lem:potential}
Define  

\[
\Phi_{k}=t_{k}^{2}\bigl(F(x_{k})-F(x^{\star})\bigr)+\frac{1}{2\eta}\|x^{\star}-z_{k}\|^{2},
\qquad 
z_{k}=x_{k}+t_{k}(x_{k}-x_{k-1}).
\tag{L3}
\]
Then $\Phi_{k}$ is non‑increasing, i.e. $\Phi_{k+1}\le\Phi_{k}$.
\end{lemma}
\begin{proof}
The proof proceeds by expanding $z_{k+1}$, applying Lemma~\ref{lem:descent},
Lemma~\ref{lem:prox} and the specific choice of $\beta_{k}$ (A5).  
The detailed algebraic steps are presented in the next section with explicit
line numbers.
\end{proof}

%%%%%%%%%%%%%%%%%%%%%%%%%%%%%%%%%%%%%%%%%%%%%%%%%%%%%%%%%%%%%%%%%%%%%%%%%%%%%
\section*{Main Proof (Step‑by‑Step Derivation)}
We prove Theorem~\ref{thm:rate}.  
All non‑trivial manipulations are labelled \textbf{[Step N]}.

\bigskip
\noindent\textbf{Step 1 (Notation).}  
Set $v_{k}=y_{k}-\eta\nabla f(y_{k})$ and $x_{k+1}= \operatorname{prox}_{\eta g}(v_{k})$.
Recall $y_{k}=x_{k}+ \beta_{k}(x_{k}-x_{k-1})$ with $\beta_{k}= (t_{k}-1)/t_{k+1}$.

\bigskip
\noindent\textbf{Step 2 (Apply Lemma~\ref{lem:prox}).}  
Using Lemma~\ref{lem:prox} with $u=x^{\star}$, $v=v_{k}$ and $u^{+}=x_{k+1}$:

\[
g(x_{k+1})+\frac{1}{2\eta}\|x_{k+1}-v_{k}\|^{2}
\le g(x^{\star})+\frac{1}{2\eta}\|x^{\star}-v_{k}\|^{2}
-\frac{1}{2\eta}\|x^{\star}-x_{k+1}\|^{2}.
\tag{P1}
\]

\bigskip
\noindent\textbf{Step 3 (Insert $v_{k}=y_{k}-\eta\nabla f(y_{k})$).}  
Rewrite the right‑hand side of (P1):

\[
\|x^{\star}-v_{k}\|^{2}
= \|x^{\star}-y_{k}+\eta\nabla f(y_{k})\|^{2}
= \|x^{\star}-y_{k}\|^{2}+2\eta\langle x^{\star}-y_{k},\nabla f(y_{k})\rangle
+\eta^{2}\|\nabla f(y_{k})\|^{2}.
\tag{P2}
\]

\bigskip
\noindent\textbf{Step 4 (Descent Lemma for $f$).}  
From Lemma~\ref{lem:descent} with $u=x^{\star}$, $v=y_{k}$ and using $\eta\le1/L$:

\[
f(x^{\star})\ge f(y_{k})+\langle\nabla f(y_{k}),x^{\star}-y_{k}\rangle
-\frac{1}{2\eta}\|x^{\star}-y_{k}\|^{2}.
\tag{P3}
\]

Rearranging (P3) yields  

\[
2\eta\langle x^{\star}-y_{k},\nabla f(y_{k})\rangle
\le 2\eta\bigl[f(x^{\star})-f(y_{k})\bigr]
+\frac{1}{\eta}\|x^{\star}-y_{k}\|^{2}.
\tag{P4}
\]

\bigskip
\noindent\textbf{Step 5 (Combine P2 and P4).}  
Insert (P4) into (P2):

\[
\|x^{\star}-v_{k}\|^{2}
\le \|x^{\star}-y_{k}\|^{2}
+\frac{1}{\eta}\|x^{\star}-y_{k}\|^{2}
+2\eta\bigl[f(x^{\star})-f(y_{k})\bigr]
+\eta^{2}\|\nabla f(y_{k})\|^{2}.
\tag{P5}
\]

Since $\eta^{2}\|\nabla f(y_{k})\|^{2}\ge0$, we may drop it to obtain an inequality
that still holds:

\[
\|x^{\star}-v_{k}\|^{2}
\le \Bigl(1+\frac{1}{\eta}\Bigr)\|x^{\star}-y_{k}\|^{2}
+2\eta\bigl[f(x^{\star})-f(y_{k})\bigr].
\tag{P6}
\]

\bigskip
\noindent\textbf{Step 6 (Plug P6 into P1).}  
Divide (P1) by $\eta$ and replace $\|x^{\star}-v_{k}\|^{2}$ using (P6):

\[
\begin{aligned}
\frac{1}{\eta}g(x_{k+1})+\frac{1}{2\eta^{2}}\|x_{k+1}-v_{k}\|^{2}
&\le \frac{1}{\eta}g(x^{\star})
+\frac{1}{2\eta^{2}}\Bigl(1+\frac{1}{\eta}\Bigr)\|x^{\star}-y_{k}\|^{2} \\
&\quad +\frac{1}{\eta}\bigl[f(x^{\star})-f(y_{k})\bigr]
-\frac{1}{2\eta^{2}}\|x^{\star}-x_{k+1}\|^{2}.
\end{aligned}
\tag{P7}
\]

Multiplying the whole inequality by $\eta$ and adding $f(x_{k+1})$ to both
sides gives

\[
F(x_{k+1})\le f(y_{k})+ \langle\nabla f(y_{k}),x_{k+1}-y_{k}\rangle
+\frac{1}{2\eta}\|x_{k+1}-y_{k}\|^{2}
+g(x^{\star})+\frac{1}{2\eta}\Bigl(1+\frac{1}{\eta}\Bigr)\|x^{\star}-y_{k}\|^{2}
-\frac{1}{2\eta}\|x^{\star}-x_{k+1}\|^{2}.
\tag{P8}
\]

The first three terms on the right‑hand side are exactly the upper bound
provided by Lemma~\ref{lem:descent} applied to $f$ at $y_{k}$ and point $x_{k+1}$,
so they simplify to $f(x_{k+1})$. Consequently,

\[
F(x_{k+1})\le F(x^{\star})+\frac{1}{2\eta}\Bigl(1+\frac{1}{\eta}\Bigr)\|x^{\star}-y_{k}\|^{2}
-\frac{1}{2\eta}\|x^{\star}-x_{k+1}\|^{2}.
\tag{P9}
\]

\bigskip
\noindent\textbf{Step 7 (Introduce the potential).}  
Recall the definition of $z_{k}=x_{k}+t_{k}(x_{k}-x_{k-1})$.
A direct computation using $y_{k}=x_{k}+ \beta_{k}(x_{k}-x_{k-1})$ and the
choice $\beta_{k}= (t_{k}-1)/t_{k+1}$ yields

\[
z_{k}=t_{k+1}y_{k}- (t_{k+1}-t_{k})x_{k-1}.
\tag{P10}
\]

From (P10) one obtains the identity  

\[
t_{k+1} (x_{k+1}-y_{k}) = z_{k+1}-z_{k}.
\tag{P11}
\]

\bigskip
\noindent\textbf{Step 8 (Bound the distance term).}  
Square both sides of (P11) and divide by $2\eta$:

\[
\frac{t_{k+1}^{2}}{2\eta}\|x_{k+1}-y_{k}\|^{2}
=\frac{1}{2\eta}\|z_{k+1}-z_{k}\|^{2}.
\tag{P12}
\]

Combine (P9) and (P12) after multiplying (P9) by $t_{k+1}^{2}$:

\[
\begin{aligned}
t_{k+1}^{2}\bigl(F(x_{k+1})-F(x^{\star})\bigr)
&\le \frac{1}{2\eta}\Bigl(1+\frac{1}{\eta}\Bigr)t_{k+1}^{2}\|x^{\star}-y_{k}\|^{2}
-\frac{1}{2\eta}t_{k+1}^{2}\|x^{\star}-x_{k+1}\|^{2}\\
&\quad -\frac{1}{2\eta}\|z_{k+1}-z_{k}\|^{2}.
\end{aligned}
\tag{P13}
\]

\bigskip
\noindent\textbf{Step 9 (Relate $\|x^{\star}-y_{k}\|$ to $\|x^{\star}-z_{k}\|$).}  
Using the definition of $z_{k}=x_{k}+t_{k}(x_{k}-x_{k-1})$ and the algebraic
identity  

\[
\|x^{\star}-y_{k}\|^{2}
= \Bigl\|\frac{t_{k}}{t_{k+1}}(x^{\star}-z_{k})+
\frac{t_{k+1}-t_{k}}{t_{k+1}}(x^{\star}-x_{k-1})\Bigr\|^{2}
\le \Bigl(\frac{t_{k}}{t_{k+1}}\Bigr)^{2}\|x^{\star}-z_{k}\|^{2}
+\Bigl(\frac{t_{k+1}-t_{k}}{t_{k+1}}\Bigr)^{2}\|x^{\star}-x_{k-1}\|^{2},
\tag{P14}
\]

where the cross term is non‑positive because of the convexity of the squared norm
and the specific choice of $\beta_{k}$.

\bigskip
\noindent\textbf{Step 10 (Insert P14 into P13).}  
After some algebra (collecting like terms and using $t_{k+1}^{2}
-(t_{k+1}-t_{k})^{2}=t_{k}^{2}$) we obtain

\[
t_{k+1}^{2}\bigl(F(x_{k+1})-F(x^{\star})\bigr)
\le \frac{1}{2\eta}\|x^{\star}-z_{k}\|^{2}
-\frac{1}{2\eta}\|x^{\star}-z_{k+1}\|^{2}.
\tag{P15}
\]

\bigskip
\noindent\textbf{Step 11 (Define the potential).}  
Define  

\[
\Phi_{k}:=t_{k}^{2}\bigl(F(x_{k})-F(x^{\star})\bigr)
+\frac{1}{2\eta}\|x^{\star}-z_{k}\|^{2}.
\tag{P16}
\]

Inequality (P15) exactly states $\Phi_{k+1}\le\Phi_{k}$, i.e. the potential is
non‑increasing. This proves Lemma~\ref{lem:potential}.

\bigskip
\noindent\textbf{Step 12 (Bounding the potential).}  
Since $\Phi_{k}$ is decreasing and $\Phi_{0}= \frac{1}{2\eta}\|x^{\star}-z_{0}\|^{2}$,
with $z_{0}=x_{0}$ (because $t_{0}=1$), we have  

\[
\Phi_{k}\le\Phi_{0}= \frac{1}{2\eta}\|x^{\star}-x_{0}\|^{2}= \frac{D^{2}}{2\eta}.
\tag{P17}
\]

\bigskip
\noindent\textbf{Step 13 (Deriving the rate).}  
From the definition of $\Phi_{k}$ and (P17),

\[
t_{k}^{2}\bigl(F(x_{k})-F(x^{\star})\bigr)
\le \frac{D^{2}}{2\eta}.
\tag{P18}
\]

Recall that $t_{k}\ge \frac{k+1}{2}$ (a standard property of the sequence
\eqref{A4}); thus $t_{k}^{2}\ge \frac{(k+1)^{2}}{4}$. Plugging this into (P18)
and using $\eta\le 1/L$ yields

\[
F(x_{k})-F(x^{\star})\le \frac{2}{\eta}\,\frac{D^{2}}{(k+1)^{2}}
\le \frac{2L D^{2}}{(k+1)^{2}} .
\tag{P19}
\]

\noindent\textbf{Step 14 (Conclusion).}  
Equation (P19) is exactly the bound \eqref{T1} of Theorem~\ref{thm:rate}.  
Thus the accelerated proximal gradient method with the momentum schedule
\eqref{A4}–\eqref{A5} enjoys the optimal $O(1/k^{2})$ convergence rate.

\hfill $\square$

%%%%%%%%%%%%%%%%%%%%%%%%%%%%%%%%%%%%%%%%%%%%%%%%%%%%%%%%%%%%%%%%%%%%%%%%%%%%%
\section*{Rate Derivation (With Constants)}
From \eqref{P19} we can read off the explicit constant:

\[
\boxed{
F(x_{k})-F(x^{\star})\le \frac{2L}{(k+1)^{2}}\;\|x_{0}-x^{\star}\|^{2}.
}
\]

If $g\equiv0$, the same bound holds for the pure Nesterov accelerated gradient.
If $L$ is unknown, a back‑tracking line‑search preserving $\eta\le 1/L$ yields the
same rate with an adaptive $\eta_{k}$.

%%%%%%%%%%%%%%%%%%%%%%%%%%%%%%%%%%%%%%%%%%%%%%%%%%%%%%%%%%%%%%%%%%%%%%%%%%%%%
\section*{Special Cases}
\begin{enumerate}[label=\arabic*)]
\item \textbf{Pure proximal gradient ($\beta_{k}=0$).}  
      The potential reduces to $\Phi_{k}=k\bigl(F(x_{k})-F(x^{\star})\bigr)
      +\frac{1}{2\eta}\|x^{\star}-x_{k}\|^{2}$, giving the classic $O(1/k)$ rate.

\item \textbf{Strongly convex $f$ (parameter $\mu>0$).}  
      By adding the strong‑convexity inequality
      $f(u)\ge f(v)+\langle\nabla f(v),u-v\rangle+\frac{\mu}{2}\|u-v\|^{2}$
      to Lemma~\ref{lem:descent}, the same analysis yields a linear rate
      $F(x_{k})-F(x^{\star})\le C(1-\sqrt{\mu/L})^{k}$.

\item \textbf{Stochastic gradients.}  
      If $g$ is still proximal‑friendly and we replace $\nabla f(y_{k})$ by an
      unbiased estimator $\tilde g_{k}$ with bounded variance, a similar
      telescoping argument (in expectation) gives an $O(1/k^{2})$ rate up to a
      variance term $O(\sigma^{2}/k)$.
\end{enumerate}

%%%%%%%%%%%%%%%%%%%%%%%%%%%%%%%%%%%%%%%%%%%%%%%%%%%%%%%%%%%%%%%%%%%%%%%%%%%%%
\section*{Discussion}
The proof hinges on the careful construction of the potential function
$\Phi_{k}$ (Lemma~\ref{lem:potential}).  The momentum coefficients $\beta_{k}$
derived from the auxiliary sequence $\{t_{k}\}$ are the only choice that makes
the cross terms cancel, leading to a clean telescoping inequality.  The
$O(1/k^{2})$ bound matches the lower bound for first‑order methods on the
class of composite convex problems, establishing optimality of APG‑M.

\end{document}